\section{The ordinary weight-one quotient at inert coefficient primes}
\label{sec:cubic-inert-coefficient}

The use of a CM family in AC-K29-ray can be replaced, at an
inert coefficient prime, by the full ordinary cohomology at each
auxiliary level.  The specializations in this section are actual
weight-one Hecke quotients.  No positive-weight ordinary CM family
at an inert prime is asserted.

\begin{proposition}[AC-K29-inert-ray: an integral inert-prime ray quotient]
\label{AC-K29-inert-ray}
Let $K=\mathbb Q(\omega)$, let $\chi$ be the cubic character
of conductor $\mathfrak f=(269)$, and let $\ell\ge5$ be inert
in $K$, with $\ell\ne269$.  Fix $a\in\{1,2\}$ and a finite
extension $E/\mathbb Q_\ell$ containing the values of $\chi$,
with integers $\mathcal O$.  If $\mathfrak n$ is prime to a
fixed finite set containing $6\ell\mathfrak f$, put
\[
 \Gamma_{\mathfrak n}=\operatorname{Gal}(K(\mathfrak n)/K),
 \qquad A_{\mathfrak n}=\mathcal O[\Gamma_{\mathfrak n}],
                                                        \tag{CIC.1}
\]
where $K(\mathfrak n)$ is the maximal abelian $\ell$-extension
of conductor dividing $\mathfrak n$.  Use the geometric Artin
normalization of LLZ, \S5.2.  For either sign $s\in\{+,-\}$
there is an actual quotient $\mathcal H_{\mathfrak n,s}$ of
the full ordinary cohomology of tame level
$N_{\mathfrak n}=|D_K|\operatorname N(\mathfrak f\mathfrak n)$,
specialized at weight coordinate $k'=-1$, with the following
properties:
\[
 \mathcal H_{\mathfrak n,s}\simeq
 \operatorname{Ind}_{G_{K(\mathfrak n)}}^{G_{\mathbb Q}}
                  \mathcal O(\chi^{-a})
 \quad\text{as }A_{\mathfrak n}[G_{\mathbb Q}]\text{-modules}.
                                                        \tag{CIC.2}
\]
Both sides are free of rank two over $A_{\mathfrak n}$.
The quotient has an $\mathcal O$-saturated kernel in the
localized ordinary module.
The isomorphisms can be chosen simultaneously so that the
actual Hecke degeneracy norm maps are the natural quotient
maps of the induced modules.  All these maps are integral;
in particular their normalization introduces no denominator
growing with $\mathfrak n$.

Consequently, for every fixed ordinary newform $f_0$ of
weight $2e+1$, $e\ge1$, level prime to $\ell$, and nebentypus
$\eta_K$, the universal ordinary Beilinson--Flach classes
give two ray Euler systems for
\[
 U_a=(M_{f_0}(e)|_{G_K}\chi^a)^*(1).
                                                        \tag{CIC.3}
\]
They have one lattice denominator depending only on
$f_0,\ell$ and the fixed tame data.  They are Bloch--Kato
away from $\ell$, and satisfy the degree-two norm relation
(CER.3).  At $\mathfrak n=1$ their quotient maps are exactly
the two ordinary Hida quotient maps for
$\theta_{\chi^a}$ with $U_\ell$ eigenvalues $+1,-1$.
The same construction permits a fixed finite cyclotomic
character and fixed tame test-vector operators as in
AC-K29-ramified and AC-K29-primitive-data.
\end{proposition}

\begin{proof}
\proofstep{1. Actual weight-one Hecke points.}
For a character $\xi$ of $\Gamma_{\mathfrak n}$ the form
$\theta_{\chi^a\xi}$ is cuspidal: the two inducing characters
remain distinct after reduction, since their ratio retains
the nontrivial cubic character and $\ell\ne3$.
It is unramified at $\ell$.  Its local Hecke polynomial is
\[
        X^2-\xi((\ell)),                                \tag{CIC.4}
\]
because $\ell$ is inert and $\chi((\ell))=1$.
The latter equality follows from anticyclotomicity: a
rational principal ideal is conjugation-fixed, so its odd
order character value equals its inverse and is one.
Squaring is an automorphism of the finite $\ell$-group
$\Gamma_{\mathfrak n}$.  Hence
$u_{\mathfrak n}=[(\ell)]^{1/2}\in A_{\mathfrak n}^\times$
is defined without adjoining coefficients.  The two roots
in (CIC.4) are $s\xi(u_{\mathfrak n})$ and
$-s\xi(u_{\mathfrak n})$; their residual values are $s,-s$.
They are distinct in characteristic $\ell\ge5$.

The source input here is precisely KLZ,
arXiv:1503.02888v6, \S7.6, Theorem~7.6.1 and the two
following paragraphs, printed p.~39.  It puts each
prime-to-$\ell$ classical weight-one stabilization in a
Hida family.  Its non-Eisenstein, residually distinguished
case identifies the actual specialized ordinary cohomology,
free of rank two, with the Artin representation.  Thus it
applies to every form just constructed.  At an auxiliary
prime dividing the chosen level but not the primitive
conductor, use its ordinary oldform, stabilized or depleted
at that auxiliary prime.  These operations are the usual
finite degeneracy maps away from $\ell$ and commute with
the ordinary projector.  They also apply to the Hida family
through the primitive form.

Write $\mathbb T_{\mathfrak n,s}$ for the completed full
ordinary Hecke algebra at the resulting residual ideal.
It is important that this is the full local algebra, not
an arbitrarily chosen characteristic-zero branch.  The
characters at all $\xi$ assemble to a continuous map
\[
 \varphi_{\mathfrak n,s}:\mathbb T_{\mathfrak n,s}
                  \longrightarrow A_{\mathfrak n},
 \qquad k'\longmapsto-1,\quad U_\ell\longmapsto s u_{\mathfrak n}.
                                                        \tag{CIC.5}
\]
Here $k'=-1$ means the character $z\mapsto z^{-1}$ on the
weight algebra in KLZ's inverse-diamond convention.
For completeness, the assembly in (CIC.5) is integral.
Away from the level its Hecke coefficient is the sum of
$\chi^a(\mathfrak q)[\mathfrak q]$ over ideals of norm $q$;
at the level use the same sum with the ideals required to
be prime to $\mathfrak f\mathfrak n$.  Diamond values and
$s u_{\mathfrak n}$ likewise belong to $A_{\mathfrak n}$.
Any relation in the Hecke algebra vanishes after evaluating
at every $\xi$, by the actual eigenpoints above, and hence
vanishes in the reduced algebra $A_{\mathfrak n}[1/\ell]$.
It therefore vanishes in $A_{\mathfrak n}$, which is
$\mathcal O$-torsion-free.  Taking completions preserves
the map: the common residual character sends the maximal
ideal into the maximal ideal of the complete finite local
ring $A_{\mathfrak n}$.  This proves (CIC.5), rather than
assuming a family of CM forms of higher weight.

\proofstep{2. Freeness and the integral induced representation.}
The residual representation is
$\operatorname{Ind}_K^{\mathbb Q}\overline\chi^{\,-a}$,
which is absolutely irreducible.  Its restriction at
$\ell$ has the two distinct unramified characters with
Frobenius values $+1,-1$.  LLZ (2015), Theorem~4.3.4,
printed pp.~14--15, or KLZ Theorem~7.2.8, printed p.~37,
therefore says that the localized full ordinary module
$\mathcal M_{\mathfrak n,s}$ is free of rank two over
$\mathbb T_{\mathfrak n,s}$.  Set
\[
 \mathcal H_{\mathfrak n,s}
       =A_{\mathfrak n}\otimes_{\varphi_{\mathfrak n,s}}
                       \mathcal M_{\mathfrak n,s}.
                                                        \tag{CIC.6}
\]
It is free of rank two over $A_{\mathfrak n}$, with its
actual integral Galois action.  After inverting $\ell$,
evaluation at every $\xi$ and KLZ \S7.6 identify it with
the corresponding induced Artin representation.  In LLZ's
convention this is the generic fibre of (CIC.2).

We check explicitly that no growing lattice index is hidden
in this generic identification.  Choose $h\in G_K$ whose
cubic character is $\omega$ and whose image in the maximal
abelian pro-$\ell$ quotient of $G_K$ is trivial.
Such an element exists: the finite cubic quotient and a
pro-$\ell$ quotient have no common nontrivial quotient,
so their joint image is their product.  The same statement
holds in the inverse limit of the finite pro-$\ell$
quotients.

This separator also proves that (CIC.6) is a literal
quotient, not just a scalar extension.  Write
$\vartheta(g)\in A_{\mathfrak n}^\times$ for the
universal ray character, and set
$t_g=\varphi_{\mathfrak n,s}(\operatorname{tr}\rho(g))$,
where $\rho$ is the actual free rank-two ordinary
representation over $\mathbb T_{\mathfrak n,s}$.
The generic identification just proved gives, for
$g\in G_K$,
\[
 \vartheta(g)^{-1}
  =\frac{t_{hg}-\omega^a t_g}
       {(\omega^{-a}-\omega^a)\chi^{-a}(g)}.             \tag{CIC.6a}
\]
The equality is integral by torsion-freeness.  Its
denominator is an $\mathcal O$-unit.  Thus every group
element of $\Gamma_{\mathfrak n}$ belongs to the Hecke
image, so (CIC.5) is surjective.  Since (CIC.6) is free
over the $\mathcal O$-torsion-free ring $A_{\mathfrak n}$,
the kernel of the resulting quotient of
$\mathcal M_{\mathfrak n,s}$ is $\mathcal O$-saturated.

Both eigenvalues of $h$ on (CIC.6) are the
fixed values $\omega^{-a},\omega^a$.  Their difference is
a unit in $\mathcal O$.  The two polynomial projectors in
$h$ consequently split (CIC.6) into projective rank-one
$A_{\mathfrak n}$-modules, which are free because this
ring is local.  On the generic fibre every $g\in G_K$
commutes with these projectors.  The same holds integrally
by torsion-freeness, so $G_K$ preserves the two lines.
Their characters are the two prescribed universal ray
characters, since this is true after every $\xi$ and
character evaluations separate $A_{\mathfrak n}[1/\ell]$.
Complex conjugation exchanges the lines and squares to
one; its two exchange entries are units with product one.
Choosing a basis on the first line and its conjugate on
the second gives exactly the standard induced module in
(CIC.2).  This proves an integral isomorphism.  The same
argument shows that its $A_{\mathfrak n}$-linear Galois
endomorphism ring is $A_{\mathfrak n}$: $h$ forces an
endomorphism to be diagonal, and complex conjugation
forces its two entries to agree.

\proofstep{3. The actual norm isomorphisms.}
We use LLZ Theorems~4.2.8--4.2.9, Corollary~4.2.10,
and Proposition~4.3.5, printed pp.~12--15.  The last
proposition expressly proves the same stabilization and
depletion isomorphisms at level $N\ell^r$ for a
non-Eisenstein, $\ell$-distinguished ideal.  The auxiliary
degeneracy maps commute with $U_\ell$ and with the
transition maps in $r$.
They give the corresponding isomorphisms for the full
ordinary modules before specialization.  One can verify
this last passage without exchanging an arbitrary tensor
product with an inverse limit: form the maps over the
finite ordinary Hecke algebras, use ordinary control to
reduce them modulo a weight-control ideal at a finite
$\ell$-power level, and apply Proposition~4.3.5 there.
The source and target are free rank-two modules over the
same completed Hecke algebra (or its quadratic
stabilization algebra).  Their reduced maps are
surjective; Nakayama makes the maps surjective, and their
determinants are units, so they are isomorphisms.  The
control ideal is contained in the maximal ideal.  This
argument is performed at an ordinary weight-two finite
level with the appropriate finite diamond character;
it asserts no algebraic control theorem of degree $-1$.
Now base change these isomorphisms by (CIC.5).

Here is also an explicit check of their scalars.  Write
$\Theta_{\mathfrak n}(\mathfrak a)
=\chi^a(\mathfrak a)[\mathfrak a]$ for a prime-to-modulus
ideal.  Denote by $D_q$ the covariant determinant operator
in the polynomial $X^2-T_qX+D_q$; in LLZ's finite-level
notation $D_q=q\langle q\rangle$.
At the weight-one quotient its value is the determinant
of the weight-one Hecke polynomial, not its weight-two
value.  Thus for a new split prime
$q=\mathfrak q\overline{\mathfrak q}$ it is
$\Theta_{\mathfrak n}(\mathfrak q)
 \Theta_{\mathfrak n}(\overline{\mathfrak q})$.
Adjoining $\mathfrak q$ retains the other Hecke root,
and the norm map is
\[
       (\operatorname{pr}_1)_*
       -\frac{\Theta_{\mathfrak n}(\mathfrak q)}q
                      (\operatorname{pr}_2)_* .         \tag{CIC.7}
\]
For a new inert auxiliary ideal $\mathfrak q=(q)$,
$T_q$ specializes to zero and $D_q$ to
$-\Theta_{\mathfrak n}((q))$.  Corollary~4.2.10 therefore
specializes to
\[
       (\operatorname{pr}_1)_*
       -\frac{\Theta_{\mathfrak n}((q))}{q^2}
                      (\operatorname{pr}_3)_* .         \tag{CIC.8}
\]
The case where the conjugate prime is already present is
the depletion map of Theorem~4.2.9.  Repeated-prime maps
are the first degeneracy map; their extra degrees are
powers of $q$ and hence units at $\ell$.
These are precisely the maps in LLZ Definition~3.3.1 and
its Hecke-equivariance calculation, with the weight-one
determinant just specified.  Their composites commute,
by the same degeneracy identities.  No division by
$\#\Gamma_{\mathfrak n}$ occurs.

Their integral induced identifications can now be chosen
compatibly by the full patching argument of LLZ
Corollary~5.2.6, printed p.~18.  Indeed, over a quotient
$A_{\mathfrak n'}\to A_{\mathfrak n}$ two choices differ
by a unit of $A_{\mathfrak n}$, by the endomorphism
calculation in Step~2.  A unit lifts to a unit of the
local ring $A_{\mathfrak n'}$.  Correct the larger
identification by that lift.  Enumerate the prime ideals,
choose an increasing cofinal sequence of moduli, and do
this successively; the commuting norm diagrams induce
all intermediate identifications.  This proves the
simultaneous assertion in (CIC.2).

For clarity, this construction also covers $\ell=5$.
The Hecke level contains $\mathfrak f$, but (CIC.5) maps
directly to the quotient ray group $\Gamma_{\mathfrak n}$.
The freeness and all norm isomorphisms survive this base
change.  Thus one need not average over any growing ray
group.  If instead one first works over the full ray
group of $\mathfrak f\mathfrak n$, its fixed tame
$5$-factor can be removed with the fixed idempotent
denominator $5$: class number one, the ray exact
sequence, and $5\nmid6$ give the product decomposition
of the $5$-primary ray group at coprime moduli.  Either
description has a denominator bounded independently of
$\mathfrak n$.

\proofstep{4. The actual Euler classes.}
Start with the universal integral ordinary class of KLZ
Theorem~6.3.4.  Apply the fixed first-factor moment and
ordinary $f_0$ quotient exactly as in AC-K29-ray,
Step~2; its single lattice denominator is independent of
the second level.  Apply (CIC.6), the compatible
identification (CIC.2), and Shapiro.  Finally specialize
the cyclotomic character to $j=e-1$.  The coefficient is
exactly (CIC.3).  All CM-side maps were integral.

The degree-two norm calculation is the universal Hecke
calculation of KLZ Theorems~5.3.1 and~5.5.1 and LLZ
Theorem~3.5.1, with the asymmetric maps described in
AC-K29-ray, Step~4.  It applies before weight
specialization.  The determinant in that calculation is
the actual $D_q$ above, so (CIC.7)--(CIC.8) give the
same degree-two polynomial after specialization.  In
particular the inert auxiliary case uses the identity
\[
 (1-\alpha zX)(1+\alpha zX)
 (1-\beta zX)(1+\beta zX)
 =(1-\alpha^2z^2X^2)(1-\beta^2z^2X^2),                 \tag{CIC.9}
\]
where $z^2=\Theta_{\mathfrak n}((q))$; this is the
Frobenius-square polynomial over $K$.  The coefficient
twist gives $R_a=M_{f_0}(e)\chi^a$, as in (CER.3).
Thus the norm relation is for $U_a=R_a^*(1)$ with
exactly the stated Frobenius convention.

The cyclotomic universal-norm argument gives the
Bloch--Kato local condition away from $\ell$, as in
AC-K29-ray, Step~5.  At the base modulus the maps are
the defining actual ordinary Hida quotient maps in KLZ
\S7.6 and Definition~8.1.1.  Fixed tame test-vector
maps are applied before these quotients: the maps are
literally $\pi_{\mathfrak n}\circ T$, with $T$ of
fixed support.  They commute with the permitted new
auxiliary degeneracies.  A fixed finite cyclotomic
specialization likewise commutes with these maps and
inserts its character in the Frobenius polynomial.
The fixed first-factor and tame denominators are the
only denominators introduced.  This proves all claims.
\end{proof}

\begin{theorem}[AC-K29-inert: the inert coefficient-field vanishing]
\label{AC-K29-inert}
Let $\ell,K,\chi$ be as in AC-K29-inert-ray.  Let $f$ be
a non-CM cuspidal newform of weight $2e+1$, $e\ge1$,
nebentypus $\eta_K$, and put $V=M_f(e)$.  Assume the
actual local representation $V|_{G_{\mathbb Q_\ell}}$
is potentially crystalline and has an ordered stable
line with diagonal characters
\[
       \eta_K\delta,\quad\delta^{-1},\qquad
       \delta|_{I_\ell}=\tau\kappa^e|_{I_\ell},
       \qquad \tau^2=1,                                \tag{CIC.10}
\]
where $\tau$ is a Dirichlet character of $\ell$-power
conductor.  Assume the characteristic-zero derived-image
conclusion of AC-K29-split (in particular, this is
supplied for the actual ordinary Type-II forms treated
there).  Then
\[
        H_f^1(K,V\chi)=H_f^1(K,V\chi^2)=0.             \tag{CIC.11}
\]
This is a coefficient-field statement.  It imposes no
bound for the integral index of the resulting Euler
class and does not address the split, residually
indistinguishable CM refinements.
\end{theorem}

\begin{proof}
\proofstep{1. Good ordinary reduction by the quadratic twist.}
The proof of AC-K29-good-twist applies with the
unramified character $\eta_K$ retained on the first
line.  More explicitly, $\tau$ is trivial or the tame
quadratic character of conductor $\ell$.  The two
diagonal characters of $V\tau$ are unramified
characters times $\kappa^e,\kappa^{-e}$.
They are crystalline.  Potential crystallinity of
their extension implies crystallinity by the
restriction--corestriction argument on
$H^1(B_{\rm cris}\otimes\operatorname{Hom})$ in that
proof.  The corrected classical Saito compatibility
statement cited there then shows that the primitive
form $f_0=f\otimes\tau$ has level prime to $\ell$.
The first line of $M_{f_0}=V\tau(-e)$ is unramified
with unit Frobenius; the other Hecke root has valuation
$2e$.  Thus $f_0$ is ordinary.  Its nebentypus remains
$\eta_K$, it is non-CM, and its derived image is the
same as that of $V$.

Apply AC-K29-inert-ray to $f_0$, then specialize the
cyclotomic variable to $j=e-1+\tau$.  The two copies
of $\tau$ cancel, exactly as in (CRC.7), and the
coefficient becomes $(V\chi^a)^*(1)$.  The actual
rank-two ray norm relation and its one fixed lattice
denominator survive this specialization.

\proofstep{2. Nonzero primitive ordinary reciprocity.}
Use the fixed primitive tame data of
AC-K29-primitive-data and the ordinary scalar
reciprocity of AC-K29-primitive-reciprocity.  The
weight-one Hida realization needed in that proof is
now the actual quotient in KLZ \S7.6, supplied by
(CIC.6).  No step of the scalar interpolation requires
that its ordinary family be a CM family.  KLZ
Theorem~7.7.2 and its following sentence expressly
allow a classical weight-one second specialization;
the nonzero rank-one ordinary functional is the
canonical map of Proposition~10.1.1(1).  The actual
universal classes and their quotient maps are being
used, so no comparison with Coleman cohomology is
inferred from scalar interpolation.

At the base modulus the two roots of $g=\theta_\chi$
are $+1,-1$.  Both are units and distinct.  The
unramified interpolation and logarithm factors have
the same valuations as (CER.5a): the $f_0$ roots
have valuations $0,2e$.  Possible unit equalities
when $e=1$ are excluded by the complex absolute
value $\ell^e$ of an $f_0$ root.  Thus all displayed
Euler and logarithm factors are nonzero.
For nontrivial $\tau$, the exact ramified formula
of KLZ Theorem~8.2.8 and Loeffler's Proposition~2.10
is the one used in AC-K29-ramified.  Its Gauss sum,
powers of $\ell$, and unit-root denominator are
nonzero, also when $\alpha_g\beta_g=-1$.
The primitive complex value is
\[
        L^{\rm prim}(f_0,g,\tau^{-1},e)
                    =L^{\rm prim}(f,g,e)\ne0,
                                                        \tag{CIC.12}
\]
by the primitive nonvanishing used in
AC-K29-primitive-reciprocity.  The fixed tame
Schwartz/test-vector choice removes the shared bad
factors, including any at $269$.  Choose one
permitted smoothing integer $c$ with $c^2\ne1$
modulo $\ell$.  Its scalar at the two rational base
classes is $c^2-1$; at higher ray levels retain the
universal group-ring smoothing factor.
Consequently each of the two base classes has a
nonzero ordinary reciprocity scalar.

\proofstep{3. The inert local comparison and the rank-two bound.}
Let $W=V\otimes\operatorname{Ind}_K^{\mathbb Q}\chi$.
Locally at $\ell$ the second factor is the sum of
the two unramified Frobenius lines $+1,-1$.
The functorial ordinary Perrin--Riou specialization
in the proof of AC-K29-lower identifies each
nonzero scalar with the dual exponential of its
actual singular class.  The $F^{--}$ component
vanishes by KLZ Proposition~8.1.7.  Hence the two
classes have nonzero singular images in the two
different one-dimensional Frobenius lines.  They
span $H_s^1(\mathbb Q_\ell,W^*(1))$, of dimension two.
For nontrivial $\tau$ the same statement follows
from the ramified local dimension calculation in
AC-K29-ramified: the finite ramified scalar makes
both crystalline modules zero, the dual exponential
is an isomorphism onto the two-dimensional
$\operatorname{Fil}^0$, and the two ordinary
refinements give its two independent lines.

The local Poitou--Tate comparison (CER.7) now forces
the localization of $H_f^1(\mathbb Q,W)$ at $\ell$
to be zero.  Under Shapiro the inert place is
\emph{one} prime $\lambda$ of $K$, of local degree
two.  Thus the conclusion is
\[
 H_f^1(K,V\chi^a)
       =H^1_{\mathrm{strict}\{\lambda\}}(K,V\chi^a).
                                                        \tag{CIC.13}
\]
It is not a comparison at two fictitious split places.
One of the constructed ray Euler systems has a
nonzero base class, because its singular localization
is nonzero.  The characteristic-zero rank-two
Euler-system bound used in AC-K29-lower applies:
the fixed lattice denominator, norm relations and
away-$\ell$ local conditions are proved above,
and the derived-image assumptions were retained.
Its dual strict group is zero.  Together with
(CIC.13) this proves (CIC.11).  Conjugation exchanges
$\chi$ and $\chi^2$ and preserves $V$, giving the
same conclusion for either choice of $a$.
\end{proof}

\paragraph{Source scope and realization dictionary.}
The new integral assertion is the assembly (CIC.5),
the separator proof of (CIC.2), and its compatibility
with the actual ordinary degeneracy maps.  The
Hida freeness, existence of the weight-one
specializations, ordinary reciprocity and
characteristic-zero Euler bound are used in their
stated source forms.  These coefficient modules are
the $\ell$-adic realizations of the theta Artin
factor and the modular-form factor; the norm maps
are the actual pull/push correspondences on the
modular curves.  No new Grothendieck attribution,
Lean execution, or uniform integral regulator
estimate is asserted.
